\documentclass[onecolumn, draftclsnofoot,10pt, compsoc]{IEEEtran}
\usepackage{graphicx}
\usepackage{url}
\usepackage{setspace}

\usepackage{geometry}
\geometry{textheight=9.5in, textwidth=7in}

\usepackage{pgfgantt}
\definecolor{grey}{rgb}{0.57, 0.64, 0.69}

% 1. Fill in these details
\def \CapstoneTeamName{The Cleverly Named Team}
\def \CapstoneTeamNumber{20}
\def \GroupMemberOne{Cade Raichart}
\def \GroupMemberTwo{Daniel Kato}
\def \GroupMemberThree{Nathan Shepherd}
\def \CapstoneProjectName{ISS Barometer App }
\def \CapstoneSponsorCompany{NASA}
\def \CapstoneSponsorPerson{Don Pettit}

% 2. Uncomment the appropriate line below so that the document type works
\def \DocType{	%Problem Statement
				% Requirements Document
				%Technology Review
				Design Document
				%Progress Report
				}

\newcommand{\NameSigPair}[1]{\par
\makebox[2.75in][r]{#1} \hfil 	\makebox[3.25in]{\makebox[2.25in]{\hrulefill} \hfill		\makebox[.75in]{\hrulefill}}
\par\vspace{-12pt} \textit{\tiny\noindent
\makebox[2.75in]{} \hfil		\makebox[3.25in]{\makebox[2.25in][r]{Signature} \hfill	\makebox[.75in][r]{Date}}}}
% 3. If the document is not to be signed, uncomment the RENEWcommand below
%\renewcommand{\NameSigPair}[1]{#1}

%%%%%%%%%%%%%%%%%%%%%%%%%%%%%%%%%%%%%%%
\begin{document}
\begin{titlepage}
    \pagenumbering{gobble}
    \begin{singlespace}
    	\includegraphics[height=4cm]{coe_v_spot1}
        \hfill
        % 4. If you have a logo, use this includegraphics command to put it on the coversheet.
        \includegraphics[height=4cm]{NASA}
        \par\vspace{.2in}
        \centering
        \scshape{
            \huge CS Capstone \DocType \par
            {\large\today}\par
            \vspace{.5in}
            \textbf{\Huge\CapstoneProjectName}\par
            \vfill
            {\large Prepared for}\par
            \Huge \CapstoneSponsorCompany\par
            \vspace{5pt}
            {\Large\NameSigPair{\CapstoneSponsorPerson}\par}
            {\large Prepared by }\par
            Group\CapstoneTeamNumber\par
            % 5. comment out the line below this one if you do not wish to name your team
            %\CapstoneTeamName\par
            \vspace{5pt}
            {\Large
                \NameSigPair{\GroupMemberOne}\par
                \NameSigPair{\GroupMemberTwo}\par
                \NameSigPair{\GroupMemberThree}\par
            }
            \vspace{20pt}
        }
        \begin{abstract}
        % 6. Fill in your abstract
This Requirements Document is meant as a repeatable and detailed description of using the iPad Air 2 Barometer to create an app for displaying barometer readings.
This app will be used to measure atmospheric pressure in emergency situations aboard the International Space Station.
This document will also serve to be an agreement that both our team and our client will make, and it will be used to evaluate our product at the end of Spring term.
Within this document is the description, constraints and dependencies that will be used for this app.
The Specific requirements that must be completed are laid out within the user, hardware and software interfaces sections.
The System features entails the more in depth features of the pressure readings, recording button, graph and the settings pages.
        \end{abstract}
    \end{singlespace}
\end{titlepage}
\newpage
\pagenumbering{arabic}
\tableofcontents
% 7. uncomment this (if applicable). Consider adding a page break.
%\listoffigures
%\listoftables
\clearpage

\section{Color Palette}

\section{Introduction}

\section{Glossary of Terms}
\begin{itemize}
\item[--] csv: A file format readable by Excel and other spread sheet software
\item[--] iOS: The internal operating system of Apple products, specifically iPads
\item[--] iPad: A touch screen tablet device with barometer capabilities
\item[--] ISS: International Space Station
\item[--] MANOVACUMETER: The current mechanical air pressure gauge aboard the ISS
\item[--] mmHg: Millimeters of Mercury, the standard pressure unit
\item[--] NASA: National Aeronautics and Space Administration
\item[--] t-Res: The chart that NASA uses to calculate residual time left on station
\end{itemize}

\section{Data Storage}

\section{Language}

\section{UI Library}

\section{Graph Interface}

\section{Barometer Implementation}

\section{Organization of UI}

\section{Chart Library}

\section{Development Framework}

\section{Conclusion}

\end{document}
